\section{Evaluation}
\label{sec:eval}

\subsection{Verification Performance}
\label{sec:eval:vpf}

We used the PPV framework to verify multiple implementations of paging functions: \icode{map}, \icode{unmap}, and \icode{resolve}, that manage multi-level page tables with dynamic page table allocation. We adapted the \icode{set\_npt} and \icode{walk\_npt} functions from the SeKVM hypervisor artifact~\cite{spoq-artifact}
%~\cite{sekvm-artifact-sosp21}, 
which respectively perform page-table updates and resolve page-table entries. Note that the \icode{unmap} primitive that we proved supports unmapping a page at the page table leaf via a \icode{map} to a NULL pte. We extracted the paging functions and primitives used to implement the functions, such as the address translation helper and the memory allocator, from SeKVM into independent C files. We fed them into the LLVM-based compiler in our proof framework. We verified that these functions satisfy the array axiom.
%refine respective specifications defined over a flat array.

%The original paging function implementations in SeKVM support translating input addresses with 3- and 4-level page tables using hybrid 2 MB and 4 KB mappings on Arm-based environments.
To demonstrate versatility, we updated SeKVM's implementations and verified two sets of paging functions that support (1) a fixed single and (2) hybrid mapping sizes. We summarize the evaluated implementations in \autoref{tbl:pf-profiles}. %In brief, the paging functions that operate on 2-level page tables support 1 GB fixed mappings, 3-level page tables support either 2 MB fixed mappings, or hybrid 1GB, 2MB, and 4KB mappings, whereas those for 4-level page tables support only 4KB mappings.
Note that we exclude paging functions operating over 4-level page tables with mixed 2 MB and 4 KB mappings~\cite{sekvm-artifact-sosp21-pt-refine} from \autoref{tbl:pf-profiles}. In our experiments, verifying this implementation caused the SMT solver to time out after 24 hours. This case highlights a current limitation in our specification design, and we plan to refine it to reduce solver complexity and improve tractability in future work.

\begin{table}[t]
\centering
\caption{SMT-Verified Paging Function Implementations}
%\renewcommand{\arraystretch}{1}
\label{tbl:pf-profiles}
\small
\begin{tabular}{ |p{1.5cm}|p{6.00cm}| }
    \hline
    \textbf{Profile} & \textbf{Description} \\
    \hline 
     Single-size & \textbf{1GB-2lvl}: Create 1 GB mappings in 2-level page tables. \textbf{2MB-3lvl}: Create 2 MB mappings in 3-level page tables. \textbf{4KB-4lvl}: Create 4 KB mappings in 4-level page tables.\\
    \hline
    Hybrid-size & \textbf{2MB+4KB-3lvl}: Create both 2 MB and 4 KB mappings in 3-level page tables. \textbf{1GB+2MB+4KB-3lvl}: Create 1 GB, 2 MB, and 4 KB mappings in 3-level page tables.\\ 
    \hline 
\end{tabular}
\end{table}

For the set (1) fixed-single mapping sizes, we verified paging functions that use only 1 GB, 2 MB, and 4 KB mapping sizes, respectively. For Arm, these implementations use 2-level, 3-level, and 4-level page tables. The set (2) hybrid-mapping sizes involves enhancements to these paging functions to accommodate variable-sized mappings -- we verified two implementations that operate on 3-level page tables: paging functions that support 2 MB and 4 KB mappings, and 1 GB, 2MB, and 4 KB mappings. The results demonstrated that our SMT-based method efficiently validates correct behavior across varying address spaces and mapping configurations. In particular, we successfully verified all cases, affirming the efficiency and effectiveness of our approach across various page-mapping schemes.

We evaluated verification performance on a machine with an i7-8750H CPU and 16 GB of memory, running Ubuntu 22.04, using Z3 as the SMT solver. We present the verification performance results for the 2-level and 3-level implementations using a quantifier-based structural lifting approach in \autoref{tbl:results}. We did not present results using quantifiers for hybrid mapping sizes because they caused SMT timeouts.
%page table implementation that provide $f_{\texttt{set\_pt}}$ and $f_{\texttt{walk\_pt}}$ as the helper functions, each differing in the number of page table levels. Experiments were conducted on a machine equipped with an i7-8750H CPU, 16 GB of memory, running Ubuntu 20.04, using Z3 as the SMT solver. Each variation was verified against the properties introduced in Section~2. For the 2-level and 3-level implementations, we further compared the performance of our specification design with the lifting-based approach that includes quantifiers. We present the results in \autoref{tbl:results}.
%additional tests were conducted using liftspecifications written with quantifiers.%

For quantifier-free tests, Z3 successfully verified all implementations with fixed and hybrid mapping sizes. For single-sized quantifier-free tests, Z3 completes verification in 175 seconds. In contrast, the quantifier-based approach could only verify the 2-level implementation with fixed sizes. It resulted in a solver timeout during verification of a 3-level fixed-mapping implementation. \autoref{tbl:results} also shows the lines of specification required in our proofs. Compared to the thousands of lines required for page table verification using foundational theorem provers like Coq~\cite{spoq}, our approach requires only modest effort to verify different page table implementations, necessitating only changes to the specification to describe the well-formedness condition for a different page-table hierarchy. 
%\todo{boundary constraints of newly added variables}.
These results demonstrate the practicality and performance advantage of our quantifier-free approach.
%Across all quantifier-free tests, Z3 successfully reasoned about the proof within 15 seconds. In contrast, lifting-based verification using quantifiers resulted in solver timeouts (exceeding 24 hours) when the structure of the page table became more complicated. These results emphasize the performance advantage provided by the quantifier-free design.

%The effort required to adapt specifications across implementations was also measured. Specifically, 63 lines of specifications were modified when switching from the 2-level to the 3-level implementation, and 78 lines were changed when switching from 3-level to 4-level. This manual effort represents the only intervention required by users. Compared to the thousands of lines involved in verifying similar systems using foundational theorem provers such as Coq, the proposed approach enables efficient adaptation and verification of different page table implementations with minimal user effort.

%Quantifiers: Whether the specifications contain quantifiers. 
\yfc{\textcolor{blue}{[Review: update experimental tables] Please update the verification-performance tables to include the new recursive page-table and modular allocator results. Identify each implementation's page-table depth, supported mapping sizes, and allocator configuration. Report the recursive verification time of 1,978 seconds and the undecomposed contract check that did not finish within eight hours. For allocator modularization, distinguish page-table verification against the abstract interface from verification with the buddy allocator inlined, and report the allocator's separate verification cost. Replace ``under a minute,'' ``over a thousand queries,'' and ``several hours'' with exact measurements where available. State the hardware, solver version, timeout, and timing scope for the new experiments; do not combine results from different configurations without identifying them. Also clarify whether any new result addresses the previously reported four-level mixed-page-size timeout.}}

\begin{table}
\centering
\caption{Evaluation of three SeKVM-like page table implementations. Loc: Lines of C code. Spec: Lines of specification. Time: Verification time in seconds, N/A indicates the solver did not terminate within 24 hours.}
\label{tbl:results}
{\footnotesize
 \begin{tabular}{cccr}
 \toprule
 \footnotesize
 \textbf{Test} & \textbf{LoC (C code)} & \textbf{Spec} & \textbf{Time(s)} \\
 \midrule
 1GB-2lvl\_quantifier & 98 & 78 & 0.215\\
 \midrule
 2MB-3lvl\_quantifier & 131 & 96 & N/A\\
 
 \midrule
 1GB-2lvl & 98 & 52 & 0.221 \\
 \midrule
 2MB-3lvl & 131 & 84 & 9.994 \\
 \midrule
 4KB-4lvl & 176 & 120 & 174.659\\
 \midrule
 2MB+4KB-3lvl & 172 & 99 & 1944.671 \\
  \midrule
 1GB+2MB+4KB-3lvl & 204 & 105 & 11343.450 \\
 \bottomrule
 \end{tabular}}
\end{table}
%[SW: modify spec to describe the boundary constraints newly added variables.]
%s (numbers in parentheses indicate the line differences between the $n$-level and $(n{-}1)$-level implementations)
% 3-Level & 132 & 84(63) & 8 \\
%\midrule
%4-Level & 152 & 120(78) & 14\\

%\sw{SW: to integrate with SeKVM and Secvma, discuss how we do that. Discuss the benefits of using an SMT-based approach by comparing it against a Rocq-based one. Add performance experimental setup + evaluation results running application benchmarks on m400.}

{\color{blue}
\paragraph{Recursive page-table verification.}
We extended PPV to verify a recursive page-table implementation.
The key observation is that recursion follows the page-table hierarchy,
whose maximum depth is fixed by the supported configuration.
We can therefore unfold the recursive specification to this depth and
express it level by level, using the same quantifier-free array encoding
as for fixed-depth implementations.
This lets the solver reason about memory reads and writes without
quantifier instantiation.
To control query size, we decompose the proof obligations by visitor
and recursion level, and split each property into cases according to
the execution path.
Without this decomposition, checking the recursive contract did not
finish within eight hours.
With the decomposition, verification of the recursive implementation
completed in 1,978 seconds.
}

{\color{blue}
\paragraph{Modular allocator verification.}
We also verified fixed-depth page-table implementations against a
modular allocator interface, supporting interchangeable SeKVM linear
allocators and a buddy allocator.
The page-table specifications refer to the allocator through a small
interface describing pool well-formedness, whether a page was allocated
from a pool, and the next page to be allocated.
The same page-table specifications can thus be reused across these
allocator implementations, while allocator code is excluded from
the page-table verification queries.
With the buddy allocator inlined, verification of the four-level
page-table implementation required over a thousand queries and several
hours, even after splitting queries into cases.
Against the abstract allocator interface, page-table verification
completed in under a minute.
This time excludes verification of the allocator itself:
the buddy allocator is checked separately against its specification,
instead of being reverified within every page-table query.
}

\subsection{Bug Identification}

During the verification of paging functions, PPV successfully uncovered three bugs. These bugs were not present in the original SeKVM source code; they were introduced unintentionally during the extraction of the paging functions into independent C files for evaluation. This underscores the inherent difficulties in designing functionally correct paging functions that incorporate complex operations—such as bit-level arithmetic—and thus demonstrates the critical need for an automated verification tool.

The first bug happened when retrieving the table index for the PGD level.
%\texttt{pgd\_idx = (((addr) >> 30) \& 0x1FF)} was modified to \texttt{pgd\_idx = (((addr) >> 30) \& 0x3FF)},
We used the wrong bit mask when performing the AND operation on the input address to the paging functions. This incorrectly expands the extracted index from 9 bits to 10 bits, resulting in incorrect entry references. %This resulted in referencing a wrong level-1 page table entry.
%\todo{unsure why PGD uses 8-bits?}

\begin{listing}
\caption{An Incorrect implementation of \texttt{el2\_va}}
\label{code:el2}
\begin{minted}[
  linenos,
  frame=lines,
  framesep=2mm,
  fontsize=\footnotesize,
  funcnamehighlighting,
  autogobble,
  xleftmargin=2em,
]{cpp}
static void *el2_va(int64_t x) {
	return ((x & (0x400000UL - 1)) | 0x400000UL);
}

void set_npt(int64_t id, int64_t addr) {
        ...
        if (pmd == 0UL) {
            pmd = alloc_pmd() | PT_FLAG;
            *(int64_t *)el2_va(pmd_addr) = pmd;
        }
        pte_idx = ((addr >> 12) & 0x1FF);
        pte_addr = (pmd & ~(0xFFF)) + pte_idx*8;
        pte = *(int64_t *)el2_va(pte_addr);
        ...
}
\end{minted}
\end{listing}

The second bug originated in an address translation primitive \texttt{el2\_va}, that converts an input physical address to a virtual address within the core's address space in Arm's EL2 mode. This primitive is essential and must be invoked before accessing memory because direct access to physical memory is not possible. The incorrect implementation of \texttt{el2\_va} applied a mask that preserves only the lower 22 bits of the input. Therefore, any physical address that differed outside these 22 bits would be truncated to the same virtual address. %This caused a problem since \texttt{alloc\_pmd} (line 8, Listing~\ref{code:el2}) can return \texttt{pmd} values whose upper bits (those that are larger than 22 bits). In this range, there are distinct \texttt{pmd} addresses differ in bits above 22. These high-order bits were discarded by \texttt{el2\_va} (line 13), causing two different \texttt{pmd} addresses to be consequently translated into the same \texttt{pte\_addr}.
The issue stems from the behavior of \texttt{alloc\_pmd} (line 8, Listing \autoref{code:el2}), which can return pmd values with the upper bits (those greater than 22) set. Addresses that were distinct only by these high-order bits were being truncated by \texttt{el2\_va} (line 13). This truncation mapped two fundamentally different \texttt{pmd} addresses to an identical \texttt{pte\_addr}. As a result, when this derived address was used for page table updates (lines 12 and 13), the system could incorrectly write to and corrupt an unrelated page table.

%from a 16 GB range (24 bits)

The third bug came from not constraining the input \texttt{pte} value of \texttt{map}. In our hybrid implementation—where paging functions support both 2 MB and 4 KB mappings—we use the lowest two bits of a page-table entry to determine whether it is a last-level entry. If both bits are set to 1, the entry is treated as an intermediate-level entry, and the traversal continues. However, an incorrectly formed \texttt{pte} that should appear in the last level may also have its two lowest bits set to 1. In that case, the traversal mistakenly treats it as an intermediate entry, resulting in an erroneous walk into the next level.

In all cases, the SMT solver quickly found a counterexample within a few seconds. The counterexample included concrete values for \texttt{mem}, \texttt{id}, \texttt{addr}, and \texttt{pte}, enabling us to efficiently identify the bug by inspecting the generated counterexample. These findings present the effectiveness of our approach in uncovering even subtle errors in paging functions.

\subsection{Paging Function Verification}

%ported two distinct paging function profiles—4KB-4lvl and 
\paragraph{Integrating with Spoq's proof of SeKVM}
SeKVM~\cite{spoq-artifact} uses a layered refinement strategy~\cite{li:sekvm:usenix21} to verify functional correctness. %The verified code includes an implementation of the paging function with 2MB+4KB-4lvl~\cite{npt-walk-sekvm-spoq-artifact}.
The authors modularize the paging functions into an independent layered module. Their verification requires structural lifting and was conducted manually; they are integrated with the rest of the code from other modules, which is automatically verified by Spoq. Spoq generates a low-level specification from C code, which is then combined with a handwritten high-level specification to establish a refinement relation. Proving that this refinement relation holds for the paging function requires substantial manual effort.
%The paging function refinement proofs~\cite{set-npt-refine-sekvm-spoq-artifact}.} 
%The verification of the \textit{map} function alone requires 2K LoC Coq code~\cite{set-npt-refine-sekvm-spoq-artifact}.}
%We tested the resulting implementations on the m400 Arm servers in CloudLab~\cite{cloudlab-hw}. We confirmed that both implementations deliver correct functionality.
%SeKVM's secure core uses the former to manage its own address-space mappings, and the latter to manage S2PTs

%We use PPV to verify the paging function variants we adapted from SeKVM's codebase (as mentioned in \autoref{sec:eval:vpf}), enabling SeKVM verification to be built on the correctness guarantees established by PPV's connecting proofs, thereby eliminating costly verification overhead for alternative paging function implementations.To integrate this into our framework, we instead use an SMT solver to verify that the target functions satisfy the conditions of update correctness and non-interference. We then construct a \textit{connecting proof} demonstrating that any function satisfying these conditions inherently refines to the high-level specification. Therefore, our push-button verification framework serves as a component in Spoq and integrates into its end-to-end verification workflow.
%We demonstrate how our framework can be seamlessly integrated with the Spoq verification framework.
We use PPV to verify the paging function variants adapted from SeKVM (as discussed in \autoref{sec:eval:vpf}). This approach allows SeKVM's overarching proofs to rely directly on the correctness guarantees established by PPV, eliminating the costly overhead of manually verifying alternative implementations.

Listing~\ref{code:spoq_set} shows the simplified high-level abstract specification of \icode{set\_npt} extracted from SeKVM in Coq. Because the page table supports two granularities, 2MB and 4KB, it is modeled as two distinct flat arrays, \icode{e\_lv2pt} and \icode{e\_lv3pt}. An execution of \icode{set\_npt} is therefore represented as an update to the corresponding array. This function is further wrapped by \icode{set\_npt\_spec} to manipulate the internal data structure, \icode{RData}, which represents SeKVM's machine model. As shown in Listing~\ref{code:spoq_set_conn}, the connecting theorem for \icode{set\_npt} formalizes Property 1 introduced in \autoref{sec:impl} and establishes that if the concrete paging implementations (\icode{set} and \icode{walk}) satisfy the two array-axiom conditions, a single execution of both the concrete \icode{set} and the abstract \icode{set\_npt\_spec} will maintain the refinement relation \icode{refrel}. Specifically, this relation \icode{refrel} ensures that the result of a concrete page table walk is equivalent to a lookup in the abstract flat arrays of \icode{RData}. The complete connecting proof for both \icode{set\_npt} and \icode{walk\_npt} is detailed in \autoref{sec:cp}. Once the connecting theorem is proven, any concrete implementation verified by PPV inherently refines to Spoq's high-level specification without additional effort. Ultimately, PPV facilitates a fully automated, push-button workflow that integrates effortlessly into Spoq’s end-to-end verification methodology.

\begin{listing}
\caption{The high-level specification of \icode{set\_npt}}
\label{code:spoq_set}
\begin{lstlisting}[style=coq, escapeinside={(*@}{@*)}]
Definition set_npt (v_addr: Z) (v_level: Z) (v_pte: Z)
(npt: Stage2PT) : (option Stage2PT) :=
  let v_gfn := (v_addr / 4096) in
  let v_gfn512 := (v_gfn / 512) in
  if (v_level =? 2)
  then (
    (Some (npt.(e_lv2pt) :< 
    ((npt.(e_lv2pt)) @@ v_gfn512 == (Some (v_pte))))))
  else (
    rely ((((npt.(e_lv2pt)) @ v_gfn512) = None));
    (Some (npt.(e_lv3pt) :< 
    ((npt.(e_lv3pt)) @@ v_gfn == (Some (v_pte)))))).

Definition set_npt_spec (v_vmid: Z) (v_addr: Z)
(v_level: Z) (v_pte: Z) (st: RData) : (option RData) :=
  let npt := (((st.(shared)).(e_s2pts)) @ v_vmid) in
  when npt' == ((set_npt v_addr v_level v_pte npt));
  (Some (st.(shared).(e_s2pts) :< (((st.(shared)).(e_s2pts)) @@ v_vmid == npt'))).
\end{lstlisting}
\end{listing}
\begin{listing}
\caption{The connecting theorem of \icode{set\_npt}}
\label{code:spoq_set_conn}
\begin{lstlisting}[style=coq, escapeinside={(*@}{@*)}]
Theorem set_connecting :
forall vmid addr level pte pte' 
    mem mem' st st' walk set
  (* hypotheses *)
  (non_interference: forall addr' level',
    addr' <> addr ->
    walk vmid addr' mem' = walk vmid addr' mem
  (update_correctness: forall addr' level',
    addr' = addr ->
    walk vmid addr' mem' = Some pte /\ pte' = pte),
  (* proof goal *)
   (refrel vmid st mem) ->
   (set_npt_spec vmid addr level pte st = Some st') ->
   (set vmid addr level pte mem = Some mem') ->
   refrel vmid st' mem'.
\end{lstlisting}
\end{listing}
\paragraph{Deploying Verified Paging Function to SECvma}
In addition, we ported the verified paging function to a prior work, SECvma~\cite{secvma2024acsac}, which employs a virtualization-based approach to protect the integrity of the Linux kernel's code. The authors of SECvma built on the SeKVM~\cite{li:ieeesp2021} hypervisor and extended SeKVM's formally verified core to include code-integrity protection of the Linux host, preventing unverified code from being loaded into the integrated Linux host kernel's memory. Specifically, SECvma's trusted core utilizes Arm's nested page tables, the stage-2 page tables (S2PT), to protect Linux's memory. Unlike SeKVM, which controls mappings (i.e., unmap VM pages) from Linux's S2PT memory isolation, SECvma builds on this to protect Linux from code-injection attacks via direct kernel memory tampering and malicious kernel page table updates. The core enables S2PT when running Linux on bare-metal hardware, hardening the kernel's memory access.
%SECvma's extensions to SeKVM's core were unverified. 
%SeKVM's verified core includes two paging functions, \texttt{set\_pt} and \texttt{walk\_pt}, that provide the \texttt{map}, \texttt{unmap}, and \texttt{resolve} functionality. These functions respectively provide page table update and walk. \texttt{set\_pt} and \texttt{walk\_pt} support traversals of multi-level page tables, dynamic page table allocations, as well as both 2MB and 4KB mappings.
Since using S2PTs incurs higher TLB pressure and slows Linux performance, the authors of SECvma adapted the integrated SeKVM's paging functions to support 1-GB huge-page mappings, in addition to 2-MB and 4-KB mappings, in S2PT to optimize TLB performance. To simplify development, the authors reimplemented the \texttt{set\_pt} function using recursion~\cite{secvma-artifact}. All extensions and adaptations to SeKVM to support SECvma, including the updated paging functions, were not verified by the authors of SECvma~\cite{secvma2024acsac}. 

To facilitate future verification efforts, SECvma's authors follow a modular approach to implement the TCB, separating paging functions into independent layer modules. We verified the \texttt{set\_pt} and \texttt{walk\_pt} functions of the 1GB+2MB+4KB-3lvl used by SECvma's core. We also ported the \texttt{set\_pt} and \texttt{walk\_pt} functions from SeKVM that use a 2MB+4KB+3-level profile, which we also verified with PPV, into SECvma to replace its paging functions. Since this implementation does not support 1GB mappings, we adapted a higher-layer module from SECvma's TCB to avoid requesting 1GB mappings in the S2PT that the core manages for Linux. 
%SECVma's core, which serves as the TCB, is unverified. 
% , 

We evaluated the performance of the host Linux kernel integrated with SECvma on the m400 hardware, with four configurations: (1) mainline Linux v4.18 (\textbf{Linux-4.18}), (2) unmodified SECvma with unverified paging functions (\textbf{SECvma}), (3) SECvma with the PPV-verified paging functions (the implementation from the 1GB+2MB+4KB-3lvl-profile) (\textbf{SECvma-ppv}), and (4) SECvma with SeKVM's paging functions (\textbf{SEC\\vma-sekvm}). We used Linux v4.18 to make a fair comparison with SECvma, whose implementation is also based on Linux v4.18. We installed the Linux kernel on an HP Moonshot m400 server with an 8-core 64-bit Armv8-A 2.4~GHz Applied Micro Atlas SoC, 64~GB of RAM, a 120~GB SATA3 SSD, and a Dual-port Mellanox ConnectX-3 10GbE NIC. For client-server workloads, clients ran on another m400 machine and connected to the server via a 10 GbE. All configurations use the Linux 4.18 kernel and Ubuntu 20.04. The kernels were compiled with the same configuration.
%The bare-metal configurations use all hardware available.
%All VMs were configured with 4 virtual CPUs and 12~GB of RAM. All VMs used paravirtualized I/O. They were configured with its standard vhost virtio network and with \texttt{cache=none} for its virtual block storage devices~\cite{tuningKVM,IBMcache}.

\begin{table}[t]
\centering
\caption{Description of benchmark application}
%\renewcommand{\arraystretch}{1}
\label{tbl:macro}
\footnotesize
\begin{tabular}{ |p{1.5cm}|p{6.00cm}| }
    \hline
    \textbf{Name} & \textbf{Description} \\
    \hline 
     Hackbench & Hackbench~\cite{hackbench} with 100 process groups running in 500 loops. Measured the time taken to execute the benchmark. Lower is better.\\
    \hline
    Apache & Apache v2.4.41 server handling 100 concurrent requests from a remote ApacheBench \cite{ab} v2.3 client, serving the 41 KB index.html of the GCC 4.4.7 manual. Measured the number of requests handled by the Apache server per second. Higher is better. \\ 
    \hline 
    Memcached & Memcached v1.5.22 using the memtier benchmark\cite{memtier} v1.2.3 with its default parameters, running eight threads. Measured throughput (ops/sec). Higher is better. \\
    \hline
    MongoDB & MongoDB server v3.6.8 handling requests from a remote YCSB\cite{cooper2010benchmarking} v0.17.0 client running the default workload A. Measured throughput (ops/sec). Higher is better. \\ % with 16 concurrent threads, readcount=10000 and operationcount=500000
    \hline
\end{tabular}
\end{table}

\begin{comment}
We evaluated the benchmarks listed in \autoref{tbl:macro} on bare-metal Linux and present the results in \autoref{tbl:app:results}. SECvma enables nested paging, which incurs a performance cost compared to mainline Linux. However, SECvma mitigates this by mapping huge pages (i.e., 1MB or 2MB) in the S2PT to optimize for the m400’s small TLB~\cite{taososp21}; consequently, the observed overhead is generally marginal. The notable exception is hackbench, where frequent multi-process context switches trigger SECvma's protection overhead. This occurs because SECvma traps updates to the sensitive Arm system registers that control paging during context switches. SECvma with verified paging functions that support 1GB mappings matches the performance of the unverified implementation. It also improves upon SeKVM's paging functions that lack 1GB S2PT mappings, particularly for hackbench, which stresses memory, by 15\%. This result shows the distinct advantage of our push-button verification approach: it significantly lowers the verification barrier of paging functions. This efficiency enables systems to support and verify newer implementations, unlocking superior performance compared to prior verified implementations.
\end{comment}

We evaluated the benchmarks shown in \autoref{tbl:macro} on bare-metal Linux and present the results in \autoref{tbl:app:results}. SECvma enables nested paging, which incurs a performance cost compared to native mainline Linux due to the incurred TLB pressure. However, SECvma mitigates this by mapping huge pages (i.e., 1 GB or 2 MB) in the S2PT to optimize for the m400’s small TLB~\cite{taososp21}. SECvma with verified paging functions that support 1 GB mappings matches the performance of the unverified implementation. It improves upon SeKVM's paging functions that lack 1 GB S2PT mappings, particularly for Hackbench, which stresses memory, by 15\%. This result shows the distinct advantage of our push-button approach, which significantly lowers the barrier to verifying paging functions. This efficiency enables systems to support and verify newer implementations, unlocking superior performance. Note that the remaining Hackbench overhead in SECvma and SECvma-ppv compared to Linux stems from the fact that SECvma traps updates to the sensitive Arm system registers that control paging. Linux frequently accesses these registers during process context switches, resulting in performance overhead. %In addition to Hackbench, we tested a few network application workloads. The results suggested that incorporating verified paging functions does not affect performance.



\begin{table}[t]
\centering
\caption{SECvma performance with verified/unverified paging functions.}
\label{tbl:app:results}
\resizebox{\columnwidth}{!}{
{\normalsize
\begin{tabular}{lcccc}
\toprule
\textbf{Config} & \textbf{Hackbench} & \textbf{Apache} & \textbf{Memcached} & \textbf{Mongodb} \\
\midrule
Linux-4.18                 & 11.30 & 6,390.55 & 86,980.19 & 397.92 \\
SECvma                     & 13.47 & 6,329.37 & 85,531.15 & 397.87 \\
SECvma-ppv             & 13.55 & 6,323.13 & 85,463.02 & 398.29 \\
SECvma-sekvm               & 15.66 & 6,279.54 & 85,847.78 & 395.78 \\
\bottomrule
\end{tabular}
}}
\end{table}
